\scanpage{11}
\mspage{9}

de ``bi-isocristal'' est à ce titre une \deleted{globalisation} \inserted{généralisation} qui englobe toutes les précédentes, et qui de plus a la louable vertu d'être stable par produit tensoriel\unsure{,|.} \inserted{et passage au dual.} Cela semble donc une catégorie toute indiquée \deleted{pour recevoir} comme catégorie des valeurs pour un foncteur ``cohomologie de De Rham'' convenable, lorsqu'on se décide à travailler à isogénie près.

1.12. \deleted{\unsure{Isocri|Isocr}} \uline{Exemple} \inserted{1}: \uline{Cas où} S \uline{est le spectre d'un corps parfait} k. Alors la donnée d'un cristal de modules de type fini sur k équivaut à la donnée d'un module de type fini M sur W, la formation du motif $\underline M^{(p)}$ correspond à la formation de
\[
M^{(p)}=M\otimes_W(W,f_W),
\]
où $f_W$ est l'endomorphisme de frobenius de W, \deleted{qui} et $(W,f_W)$ est la W-algèbre dé\unsure{fi}nie par $f_W$. Par suite un structure de p-cristal sur M revient à la donnée d'un homomorphisme de W-modules $M^{(p)}\longrightarrow M$, ou si on préfère, à la donnée d'un homomorphisme $f_W$\deleted{\unsure{,|-}} semi-linéaire
\[
F_M:M\longrightarrow M.
\]
Comme l'application $x\rightsquigarrow x\otimes1$ de M dans $M\otimes_W(W,f_W)$ est bijective, $f_W$ étant un automorphisme de W, on peut considérer la bijection inverse, qui est $f_W^{-1}$-semi-linéaire. Par suite, la donnée d'une $p^{-1}$-structure sur M revient à la donnée d'un homo\unsure{morphisme|morpbisme} $f_W^{-1}$-linéaire:
\[
V_{\inserted{M}}:M\deleted{\unsure{T}}\longrightarrow M.
\]
la donnée d'un couple $(F,V)$ définit sur M une structure de bi-cristal de poids i si et seulement si on a
\[
FV=VF=p^i\operatorname{id}_M.
\]
En particulier, pour $i=1$, on retrouve la notion de \uline{module de Dieudonné}: la catégorie des cristaux de Dieudonné sur k est canoniquement équivalente à celle des modules de Dieudonné \deleted{sur k} relativement à k.

La catégorie des isocristaux de \deleted{modules} type fini sur k est équivalente à la catégorie des vectoriels de dimension finie sur le corps des fractions \inserted{K}\deleted{$\unsure{W\otimes\underline Q_p|W_{\underline Q_p}}$} de W. Donc un bi\inserted{iso}cristal (de type fini) sur k s'identifie à un tel vectoriel, muni d'un $f_K^{\deleted{\unsure{-1}}}$-endomorphisme bijectif.
\margin{10}
